\documentclass[11pt,a4paper]{article}

\usepackage[margin=1in]{geometry}
\usepackage[T1]{fontenc}
\usepackage[utf8]{inputenc}
\usepackage{lmodern}
\usepackage{microtype}
\usepackage{amsmath,amssymb,mathtools}
\usepackage{booktabs}
\usepackage{graphicx}
\usepackage{float}
\usepackage[hidelinks]{hyperref}

\graphicspath{{../demos/report_outputs/exercise_3/}{Balint/demos/report_outputs/exercise_3/}}

\title{Final Report}
\author{}
\date{}

\begin{document}

\maketitle

\section*{Exercise 3: Variational Ground-State Search With RBM NQS}

This section studies variational ground-state optimization for the transverse-field Ising model within the neural-quantum-state workflow developed in the project. The calculation combines an RBM ansatz, local Metropolis sampling, the Adam optimizer, and Renyi-2 as the entanglement diagnostic tracked alongside the energy. The discussion is organized around the architecture choice, the VMC setup, a small-system benchmark where exact energies remain available as a clean reference, a larger-system training study, and two short outlook sections on GHZ targets and entanglement-spectrum tails.

\subsection*{3/a Architecture Choice}

The variational ansatz used in this section is \texttt{RBM(alpha=2)}. This choice is motivated by the architecture comparison from the previous exercise: under the same sampler and observable workflow, the RBM family was the only one that consistently produced clearly nontrivial half-partition Renyi-2 values.

Within that family, the higher-capacity choice \texttt{RBM(alpha=2)} is adopted here. The point is not that RBMs are universally optimal, but that they provide the most useful entanglement-carrying ansatz among the architectures tested in the same implementation.

The architecture summary supports the same reading numerically. In the comparison carried over from Exercise 2, the RBM gives the largest half-partition Renyi-2 among the three model families, while the FFNN and CNN remain substantially smaller.

\begin{figure}[H]
    \centering
    \includegraphics[width=0.82\linewidth]{exercise_3_architecture_summary.png}
    \caption{Architecture-comparison summary carried into the Exercise 3 ansatz choice.}
    \label{fig:architecture-summary}
\end{figure}

\subsection*{3/b VMC Setup For The TFIM Ground-State Search}

The variational objective is the ground-state energy of the neural quantum state,
\begin{equation}
    E_\theta = \frac{\langle \psi_\theta | H | \psi_\theta \rangle}{\langle \psi_\theta | \psi_\theta \rangle}.
\end{equation}
The code path follows the same module-based structure used in the earlier exercises. The model layer provides the RBM wavefunction through the common $\log \psi_\theta(\sigma)$ interface, the sampler module generates configurations with a local Metropolis chain, the observables module provides the Renyi-2 diagnostics, the optimizer module supplies the Adam update rule, and the VMC driver performs the step-by-step parameter updates.

The Hamiltonian expectation value is evaluated through the local-energy estimator. For each sampled configuration $\sigma$, the operator helper enumerates the connected configurations $\sigma'$ generated by the TFIM action, and the local energy is written as
\begin{equation}
    E_{\mathrm{loc}}(\sigma)
    =
    \sum_{\sigma'} H_{\sigma',\sigma}\,
    \frac{\psi_\theta(\sigma')}{\psi_\theta(\sigma)}.
\end{equation}
The gradient path is the standard log-derivative covariance estimator with
\begin{equation}
    O_\theta(\sigma) = \nabla_\theta \log \psi_\theta(\sigma),
\end{equation}
so that
\begin{equation}
    \nabla_\theta E_\theta
    =
    2\,\mathrm{Re}\!\left[
        \langle O_\theta^* E_{\mathrm{loc}} \rangle
        -
        \langle O_\theta^* \rangle
        \langle E_{\mathrm{loc}} \rangle
    \right].
\end{equation}
The VMC driver uses this sampled energy-and-gradient path directly, and Adam then turns the gradient estimate into a parameter update.

Sampling is performed with \texttt{MetropolisLocal} using 16 chains and 32 discard steps per chain. Energy is logged at every training step and the half-partition Renyi-2 is logged every two steps. The small benchmark uses a 256-step training budget with 256 measurement samples, while the larger sweep uses 70 training steps with 384 measurement samples.

One important point in this exercise is the status of the entanglement diagnostic. For the small proxy systems studied here, the trained NQS state can still be reconstructed exactly, so the subsystem-size scans use exact Renyi-2 of the trained variational state. That gives cleaner figures for these system sizes. The sampled SWAP route remains the relevant general method once exact reconstruction is no longer affordable.

\subsection*{3/c Small-System Benchmarks: Energy And Renyi-2 During Training}

The exact small-system benchmark uses a periodic $4\times 1$ TFIM chain and compares three field values under the same RBM and optimizer settings: a ferromagnetic point at $h = 0.5$, the critical point at $h = 1.0$, and a paramagnetic point at $h = 1.5$. This is the cleanest place to compare optimization behavior, because exact ground-state energies and exact half-partition Renyi-2 values are still available as references while the training histories remain easy to inspect directly. In one dimension, critical TFIM states are expected to show enhanced entanglement relative to gapped phases, with logarithmic scaling at criticality and area-law behavior away from it.\footnote{See Calabrese and Cardy, ``Entanglement Entropy and Quantum Field Theory,'' \emph{J. Stat. Mech.} P06002 (2004).}

\begin{table}[H]
    \centering
    \begin{tabular}{lrrrrr}
        \toprule
        run & final energy & exact energy & energy error & final $S_2$ & exact $S_2$ \\
        \midrule
        ferromagnetic $h=0.5$ & -4.269608 & -4.271558 & 0.001950 & 0.647272 & 0.659281 \\
        critical $h=1.0$ & -5.225886 & -5.226252 & 0.000366 & 0.386575 & 0.386461 \\
        paramagnetic $h=1.5$ & -6.758767 & -6.760009 & 0.001242 & 0.165577 & 0.167198 \\
        \bottomrule
    \end{tabular}
    \caption{Small-system summary for the periodic $4\times 1$ TFIM benchmark.}
    \label{tab:small-summary}
\end{table}

\begin{figure}[H]
    \centering
    \includegraphics[width=0.88\linewidth]{exercise_3_small_energy.png}
    \caption{Energy history for the periodic $4\times 1$ TFIM benchmark.}
    \label{fig:small-energy}
\end{figure}

\begin{figure}[H]
    \centering
    \includegraphics[width=0.88\linewidth]{exercise_3_small_entropy.png}
    \caption{Half-partition Renyi-2 history for the periodic $4\times 1$ TFIM benchmark.}
    \label{fig:small-entropy}
\end{figure}

With the longer training schedule, the small benchmark is now cleanly converged across all three field values. The final energy errors are of order $10^{-3}$ throughout, and the final half-partition Renyi-2 values also track the exact benchmarks closely. The critical run remains the most accurate of the three, but the ferromagnetic and paramagnetic runs now agree with the exact reference at the same qualitative level rather than lagging far behind it.

This improved benchmark gives a much clearer calibration of the retained RBM/VMC setup. On the periodic $4\times 1$ chain, the workflow is capable of reproducing both energies and coarse entanglement diagnostics to high accuracy once the optimization budget is made large enough. That result strengthens the interpretation of the later sampled-only sections: deviations there are much more plausibly due to large-system optimization and sampling difficulty than to a basic failure of the underlying ansatz or training loop.

\subsection*{3/d Larger-System Training}

The larger-system study keeps the same RBM, sampler, and optimizer, but shifts the emphasis from exact benchmarking to sampled training behavior. The retained systems are a periodic TFIM chain of length $32$, a periodic $5\times 5$ TFIM lattice, and a periodic $5\times 5$ square-lattice $J_1$-$J_2$ model. Exact energies and exact entanglement references are intentionally disabled in this section, so the results should be read as sampled-VMC benchmarks rather than as controlled exact comparisons.

The parameter choices are now aligned with the standard literature benchmarks. For the one-dimensional chain, the two retained points are the critical value $h=1.0$ and an off-critical paramagnetic point at $h=1.5$. For the square-lattice TFIM, the two-dimensional critical field is close to $h_c/J \approx 3.04$, so the notebook uses $h=3.0$ and $h=4.0$ to probe opposite sides of the transition.\footnote{For representative estimates of the square-lattice TFIM critical point, see for example Bl\"ote and Deng, ``Cluster Monte Carlo simulation of the transverse Ising model,'' \emph{Phys. Rev. E} 66 (2002), and subsequent tensor-network benchmarks reporting $h_c/J \approx 3.04$.} For the frustrated square-lattice $J_1$-$J_2$ model, the retained ratios are $J_1/J_2=1.8$ and $2.2$, corresponding to $J_2/J_1 \approx 0.56$ and $0.45$ and therefore straddling the commonly discussed strongly frustrated window near $J_2/J_1 \sim 0.5$.\footnote{For representative studies of the square-lattice spin-$1/2$ $J_1$-$J_2$ model and the intermediate nonmagnetic regime, see Hu, Becca, Parola, and Sorella, ``Direct evidence for a gapless $Z_2$ spin liquid by frustrating N\'eel antiferromagnetism,'' \emph{Phys. Rev. B} 88 (2013), and Bishop, Farnell, and Parkinson, ``Phase Transitions in the Spin-Half $J_1$--$J_2$ Model,'' \emph{Phys. Rev. B} 58 (1998).}

Figure~\ref{fig:large-graphs} summarizes the three lattice geometries used in the larger-system study. Blue edges denote nearest-neighbor couplings in all three cases, while the frustrated square-lattice model adds red diagonal second-neighbor bonds. For visual clarity, the two square-lattice graphs are drawn with open boundaries even though the simulated Hamiltonians use periodic boundary conditions. The one-dimensional chain is shown directly with periodic boundaries because the ring geometry remains visually clear at this size.

\begin{table}[H]
    \centering
    \begin{tabular}{lrrr}
        \toprule
        run & final energy & final $S_2$ & tail energy std. \\
        \midrule
        TFIM $L=32$, $h=1.0$ & -27.779286 & 1.266810 & 0.471005 \\
        TFIM $L=32$, $h=1.5$ & -13.777831 & 0.859072 & 0.740313 \\
        TFIM $5\times 5$, $h=3.0$ & -42.979151 & 0.221975 & 1.382138 \\
        TFIM $5\times 5$, $h=4.0$ & -49.251486 & 3.054304 & 4.485543 \\
        $J_1$-$J_2$ $5\times 5$, $J_1/J_2=1.8$ & -11.999737 & 0.979289 & 0.114004 \\
        $J_1$-$J_2$ $5\times 5$, $J_1/J_2=2.2$ & -27.194943 & 0.285725 & 0.165600 \\
        \bottomrule
    \end{tabular}
    \caption{Sampled larger-system summary for the periodic TFIM and $J_1$-$J_2$ benchmarks.}
    \label{tab:large-summary}
\end{table}

\begin{figure}[H]
    \centering
    \begin{minipage}[t]{0.32\linewidth}
        \centering
        \includegraphics[width=\linewidth]{exercise_3_tfim_chain_graph.png}
    \end{minipage}\hfill
    \begin{minipage}[t]{0.32\linewidth}
        \centering
        \includegraphics[width=\linewidth]{exercise_3_tfim_square_graph.png}
    \end{minipage}\hfill
    \begin{minipage}[t]{0.32\linewidth}
        \centering
        \includegraphics[width=\linewidth]{exercise_3_j1j2_square_graph.png}
    \end{minipage}
    \caption{Lattice graphs used to visualize the larger-system benchmarks. Left: periodic TFIM chain with $L=12$, drawn as a ring. Middle: $5\times 5$ TFIM lattice drawn with open boundaries for clarity, although the simulations use periodic boundaries. Right: $5\times 5$ $J_1$-$J_2$ lattice drawn with open boundaries for clarity; blue edges denote $J_1$ nearest-neighbor couplings and red edges denote $J_2$ diagonal couplings.}
    \label{fig:large-graphs}
\end{figure}

\begin{figure}[H]
    \centering
    \includegraphics[width=0.88\linewidth]{exercise_3_large_energy.png}
    \caption{Energy history across the periodic larger-system benchmark set.}
    \label{fig:large-energy}
\end{figure}

\begin{figure}[H]
    \centering
    \includegraphics[width=0.88\linewidth]{exercise_3_large_entropy.png}
    \caption{Half-partition Renyi-2 history across the periodic larger-system benchmark set.}
    \label{fig:large-entropy}
\end{figure}

\begin{figure}[H]
    \centering
    \includegraphics[width=0.9\linewidth]{exercise_3_large_entropy_scan.png}
    \caption{Post-training Renyi-2 versus subsystem size for the periodic larger-system benchmark set.}
    \label{fig:large-entropy-scan}
\end{figure}

The large-system figures are most useful as diagnostics of stability and qualitative entanglement ordering. In the one-dimensional TFIM chain, the expected ordering is preserved: the critical run at $h=1.0$ ends with a larger half-partition Renyi-2 than the off-critical run at $h=1.5$. This is the cleanest qualitative success in the sampled-only part of the exercise.

The two-dimensional runs are more uneven. On the periodic $5\times 5$ TFIM lattice, the $h=3.0$ run near the known critical region finishes with a modest sampled half-partition Renyi-2, while the $h=4.0$ run on the paramagnetic side ends with a much larger value and much larger late-stage energy fluctuations. The corresponding tail-energy standard deviations, $1.38$ and $4.49$, show that these two-dimensional TFIM runs are not equally well converged; the very large final $S_2$ at $h=4.0$ should therefore be interpreted as a warning sign of unstable sampled optimization rather than as a physically trustworthy entanglement ordering.

The frustrated $J_1$-$J_2$ pair is more stable than the two-dimensional TFIM pair. Both retained runs stay in a visibly narrower late-stage energy band, and their tail-energy standard deviations remain below $0.17$. Within this pair, the more frustrated point $J_1/J_2=1.8$ ends with the larger sampled half-partition Renyi-2 than $J_1/J_2=2.2$, which is qualitatively consistent with the expectation that stronger frustration near the intermediate regime enhances entanglement. Even so, these remain sampled diagnostics rather than controlled phase-boundary measurements.

\subsection*{3/e GHZ Diagnostic}

The GHZ section serves as a qualitative diagnostic of how the same RBM/VMC workflow behaves on a highly nonclassical target state, now phrased as the periodic $L=32$ TFIM chain at $h=0$.

On this large sampled-only problem the retained quantity of interest is the training history itself. The present workflow does not provide a reliable exact-state fidelity for this system size, and the sampled Renyi callback on this target still produces an unphysical negative entropy estimate in the summary output. The figure is therefore included only as evidence about optimization behavior on a cat-like target, not as a completed fidelity benchmark.

\begin{figure}[H]
    \centering
    \includegraphics[width=0.82\linewidth]{exercise_3_ghz_history.png}
    \caption{Energy history for the periodic $L=32$ TFIM chain at $h=0$, shown as a qualitative GHZ-style diagnostic.}
    \label{fig:ghz-history}
\end{figure}

\subsection*{3/f Entanglement-Spectrum-Tail Appendix}

The exact $4\times 1$ benchmark makes it possible to say something more precise about tail sensitivity than was possible in the larger sampled-only systems. Table~\ref{tab:tail-summary} and Figure~\ref{fig:tail-spectrum} compare the exact half-partition spectrum across the three small-system TFIM benchmarks and summarize both the von Neumann entropy $S_1$ and the Renyi-2 entropy $S_2$.

\begin{table}[H]
    \centering
    \begin{tabular}{lrrrr}
        \toprule
        run & exact $S_1$ & exact $S_2$ & $S_1-S_2$ & tail weight after top 2 \\
        \midrule
        TFIM $4\times 1$, $h=0.5$ & 0.681084 & 0.659281 & 0.021803 & 0.000734 \\
        TFIM $4\times 1$, $h=1.0$ & 0.521869 & 0.386461 & 0.135408 & 0.004278 \\
        TFIM $4\times 1$, $h=1.5$ & 0.309274 & 0.167198 & 0.142076 & 0.005566 \\
        \bottomrule
    \end{tabular}
    \caption{Exact tail-sensitive summary for the periodic $4\times 1$ TFIM benchmark.}
    \label{tab:tail-summary}
\end{table}

\begin{figure}[H]
    \centering
    \includegraphics[width=0.88\linewidth]{exercise_3_tail_spectrum.png}
    \caption{Exact half-partition entanglement spectrum for the periodic $4\times 1$ TFIM benchmark.}
    \label{fig:tail-spectrum}
\end{figure}

This comparison supports a conservative but useful conclusion. Exercise 2 already showed that sampled SWAP estimates reproduce exact small-system $S_2$ accurately enough for coarse entanglement benchmarking, and the present Exercise 3 benchmark retains that utility. However, $S_2$ is dominated by the largest reduced-density-matrix eigenvalues and is therefore only a blunt probe of the tail. The difference $S_1-S_2$ is much larger for the $h=1.0$ and $h=1.5$ ground states than for the $h=0.5$ state, even though all three are represented by the same four Schmidt levels on this cut. The tail weights after the two leading eigenvalues show the same pattern.

The practical implication is that agreement in energy and $S_2$ is necessary but not sufficient if the goal is to claim faithful reconstruction of the smallest Schmidt weights. The current workflow is therefore well supported as a coarse entanglement benchmark, but not yet as a tail-faithfulness benchmark for larger NQS states.

\subsection*{Conclusion}

Exercise 3 carries the project from static entanglement diagnostics to an actual variational ground-state search. The selected ansatz is \texttt{RBM(alpha=2)}, chosen because the RBM family already emerged as the most useful entanglement-capable model in the shared workflow. The VMC setup then combines that ansatz with the local Metropolis sampler, the connected local-energy estimator, the log-derivative gradient path, and the Adam optimizer.

The exact $4\times 1$ benchmark now shows that the retained RBM/VMC workflow can be quantitatively accurate across the ferromagnetic, critical, and paramagnetic points once the small-system optimization budget is sufficiently long. All three runs reach milliscale energy errors, and the final half-partition Renyi-2 values agree closely with the exact references. That benchmark therefore provides a strong calibration of the underlying ansatz and training loop before moving to larger sampled-only systems.

The larger-system study confirms that the workflow remains operational on a periodic $L=32$ TFIM chain, a periodic $5\times 5$ TFIM lattice, and a periodic $5\times 5$ $J_1$-$J_2$ model, but it also shows that interpretation must remain cautious without exact references. Some qualitative trends are physically sensible, especially the larger half-partition Renyi-2 in the critical one-dimensional TFIM run than in its off-critical counterpart and the stronger entanglement at $J_1/J_2=1.8$ than at $2.2$ in the frustrated square-lattice pair. By contrast, the $5\times 5$ TFIM comparison remains visibly unstable at the present budget, with the paramagnetic $h=4.0$ run carrying both the largest sampled Renyi-2 and the largest late-stage energy variance. The large-system figures are therefore best read as sampled benchmark diagnostics rather than definitive statements about phase-level entanglement ordering.

Finally, the tail appendix sharpens the entanglement interpretation. The current workflow is supported as a coarse benchmark through energy comparisons and reliable small-system $S_2$ checks, but the exact $4\times 1$ spectrum shows that $S_2$ alone does not determine the tail of the entanglement spectrum. That distinction is important for interpreting all later large-system NQS results: strong agreement in coarse observables is valuable, but it should not be overstated as proof of full spectral fidelity.

\end{document}
